\documentclass[acmsmall,nonacm]{acmart}

\usepackage{microtype}
\usepackage{enumitem}
\usepackage{amsthm}
\usepackage{caption}
\usepackage{booktabs}


% ---------- Simple block bracketing ----------
\newenvironment{bracketlines}{%
  \par\medskip\hrule\medskip\noindent
}{%
  \par\medskip\hrule\medskip
}

% ---------- Theorem environments ----------
\newtheorem{definition}{Definition}
\newtheorem{theorem}{Theorem}
\newtheorem{lemma}{Lemma}
\newtheorem{corollary}{Corollary}
\theoremstyle{remark}
\newtheorem{remark}{Remark}
\newtheorem{example}{Example}
\newtheorem{observation}{Observation}
\newtheorem{assumption}{Assumption}

\numberwithin{definition}{section}
\numberwithin{theorem}{section}
\numberwithin{lemma}{section}
\numberwithin{corollary}{section}
\numberwithin{remark}{section}
\numberwithin{example}{section}
\numberwithin{observation}{section}
\numberwithin{assumption}{section}

% ---------- Notation ----------
\newcommand{\Hist}{\mathcal{H}}
\newcommand{\Proc}{\mathcal{P}}
% \newcommand{\Obs}{\mathsf{Obs}} % reserved for resolved-observation function names, not the outcome set
\newcommand{\Spec}{\mathsf{Spec}}
\newcommand{\Term}{\mathsf{Term}}
\newcommand{\Ord}{\preceq}
\newcommand{\cost}{\mathsf{cost}}
\newcommand{\Cost}{\mathsf{Cost}}
\newcommand{\depth}{\mathsf{depth}}
\newcommand{\sqsubseteqh}{\sqsubseteq_h}
\newcommand{\Aop}{\mathsf{A}}
\newcommand{\Depth}{\mathsf{Depth}}
\newcommand{\Resolve}{\mathsf{Resolve}}
\newcommand{\view}{\mathsf{view}}
\newcommand{\Ag}{\mathcal{A}g}
\newcommand{\CoordCost}{\mathsf{CoordCost}}
\newcommand{\CoordDepth}{\mathsf{CoordDepth}}
\newcommand{\ICost}{\mathsf{ICost}}
\newcommand{\IDepth}{\mathsf{IDepth}}
\newcommand{\Atoms}{\mathsf{Atoms}}
% \newcommand{\Out}{\mathcal{O}} % legacy; use $O$ for the set of outcomes/observations
\newcommand{\drule}{\mathrel{:\!-}}
\newcommand{\Prov}{\mathsf{Prov}}
% Certificates denote outcome constraints.
\newcommand{\den}[1]{\llbracket #1 \rrbracket} % semantics bracket

% ---------- Comments, etc. ----------
\newenvironment{commentary}{%
  \begin{quote}\itshape\small\color{gray}
}{%
  \end{quote}
}


\title{Provenance of Interpretations: From Ambiguity to Algebra}
\author{Joseph M. Hellerstein}
\affiliation{%
  \institution{University of California, Berkeley and Amazon Web Services}
  \country{USA}
}
\date{\today}

\acmConference[Conference]{Full Conference Name}{Month Year}{Location}
\acmBooktitle{Proceedings of the Conference}
\acmYear{2024}
\acmISBN{123-4-56789-012-3}


\begin{document}

\begin{abstract}
Algebraic provenance explains query results by tracing derivations under a fixed,
resolved semantics.
Many systems, however, operate under \emph{semantic ambiguity}:
logic programs with negation admit multiple accepted models, transactional
systems admit multiple consistent executions, and distributed or policy-aware
systems admit multiple outcomes compatible with the same history.
In such settings, it is not enough to explain \emph{what follows;} one must also
explain \emph{under which semantic resolution of ambiguity} it follows.

This paper introduces \emph{interpretation provenance}, a framework for making
semantic resolution explicit and explainable.
An \emph{interpretation} is a structured collection of semantic commitments that
resolves ambiguity by refining admissible outcomes.
We show that provenance becomes meaningful only \emph{after} such resolution.
% , and
% that classical algebraic provenance arises as a special case when semantics are
% already resolved.
We formalize how interpretations structure provenance prior to resolution, giving
rise to a layered, dependency-sensitive form that becomes ordinary semiring
provenance once a resolving interpretation is fixed.
We introduce \emph{layer certificates} as sufficient summaries of semantic
commitments for downstream reasoning.

The work extends the reach of formal data provenance to a range of 
settings where ambiguity is intrinsic.
To illustrate, we instantiate the framework in two domains with fundamentally different sources
of ambiguity: negation in Datalog and scheduling in transactional systems.
Across both settings, interpretation provenance clarifies which semantic
commitments were required to justify observations, and why interpretations
are required to complement classical provenance in the presence of ambiguity.
\end{abstract}

\maketitle

% ================================================================
\section{Introduction}
\label{sec:intro}

Algebraic provenance has given database theory a powerful and durable language for
explaining query results.
Declarative query languages provide well-defined semantics,
and provenance provides clean, uniform explanations of query
results—supporting why- and how-provenance, uncertainty propagation,
and optimization through a common algebraic framework
\cite{green2007provenance,suciu2011probabilistic}.
In these settings, explanation reduces to arithmetic, and this reduction has been
remarkably effective.

Practical data systems can make use of this foundation, but they operate in a
broader semantic landscape.
Transactional workloads admit multiple serializable executions.
Logic programs with negation admit multiple distinct models.
Distributed systems can yield multiple admissible orderings consistent with the same
history.
In such settings, a history determines not a single outcome, but a structured
space of multiple admissible outcomes.

Natural provenance questions often require more involved answers in this context
than classical provenance can provide.
For example, the reason that a tuple was
returned by a query may depend on which serializable execution was realized by a scheduler.
New questions arise in this context as well that have no analog in classical provenance.
Did this tuple appear as an inevitable consequence regardless of transactional decisions,
or because of a particular scheduling choice made by the system?
Was this tuple returned due to a particular choice of isolation
level---e.g., snapshot isolation versus serializability?
These familiar transactional examples are emblematic of the gap between classical data provenance and a range of
systems research on 
lineage~\cite{sigelman2010dapper,seltzer2005provenance,bates2015trustworthy}
that has lacked a formal foundation.
Answering such questions requires explaining not only how conclusions were derived,
but \emph{how semantic ambiguity was resolved}.

Prior work developed
\emph{interpretations} as formal objects that resolve semantic ambiguity by
introducing structured semantic commitments~\cite{hellerstein2026interpretation}.
That work develops a complexity theory of interpretation, characterizing the
\emph{cost} and \emph{depth} of semantic resolution.
Interpretations explain both semantic ambiguity arising within formal languages
(as in logic programming with negation) and ambiguity arising from execution
environments (as in transactional and distributed systems).

This paper shows how interpretations give rise to a new form of provenance.
We introduce \emph{interpretation provenance}, which records the semantic
commitments under which observations become well defined.
Once semantic ambiguity has been resolved by an interpretation, classical data
provenance applies directly, providing the derivational structure of results
within that interpretation.

We illustrate interpretation provenance in two domains with fundamentally different
sources of semantic ambiguity.
Logic programming with negation provides the canonical, mathematically explicit
instance of semantic ambiguity and resolution in computer science, serving as a
control experiment for interpretation provenance.
We then turn to transactional systems, where ambiguity arises from concurrency and
isolation choices, and show how interpretation provenance provides a fuller
accounting of query results than classical approaches. In the appendix we also examine
examples from distributed systems observability, a widely-used setting for system
lineage~\cite{sigelman2010dapper}. 
Here we give an initial transactional example to
illustrate the ideas:

\begin{example}[Provenance under transactional ambiguity]
  \label{ex:example}
Consider a database with a binary relation $S$ 
and an initially empty instance (EDB).
The following set of transactions is issued concurrently:
\[
\begin{array}{ll}
T_1: \textsf{insert } S(1,b) \ (x_1) \qquad &
T_2: \textsf{insert } S(2,b) \ (x_2)\\
T_3: \textsf{insert } S(3,c) \ (x_3) \qquad &
T_4: \textsf{insert } S(4,d) \ (x_4)\\
T_5: \textsf{delete } S(4,d). &
\end{array}
\]
Subsequently, a query transaction $T_Q$ evaluates $Q(y) \drule S(k,y)$.

Under serializability, the observed history admits multiple executions.
Tuples $\{(b), (c)\}$ appear in $Q$ for every serializable execution.
Tuple $(d)$ is contingent: it appears in some executions and not others, depending
on whether the delete is ordered before or after the query. Before resolving this
ambiguity, even the \emph{family} of provenance expressions for $(d)$ is ambiguous:
in one order the question is \emph{why} $d$ appears, while in another it is \emph{why-not} $d$.
Interpretation provenance makes this dependence explicit.
It records the semantic commitments under which the query result is interpreted,
together with the ordinary derivational explanation that applies once those
commitments are fixed.
Section~\ref{sec:concurrency} formalizes this example by writing down the corresponding
resolving interpretations and provenance expressions, and the example is 
worked in full in Appendix~\ref{sec:appendix-concurrency}.
\end{example}

\paragraph{Contributions.}
This paper introduces \emph{interpretation provenance}: a semantic layer that
explains conclusions in the presence of ambiguous semantics by recording the
commitments under which those conclusions become determinate.
We (i) formalize interpretations and certificates as the semantic structure
needed to recover classical algebraic provenance once ambiguity is resolved;
(ii) show, via Datalog with negation, how different semantic regimes induce
distinct interpretation and provenance structures, (iii) show, via transactional
systems, how concurrency and isolation choices give rise to interpretation-dependent
query results; and (iv) identify \emph{robustness} across interpretations and
\emph{preservation} across changes in specifications
as new classes of provenance questions enabled by this framework.


% \begin{enumerate}
% \item It identifies \emph{semantic ambiguity}—the presence of multiple admissible
% outcomes for a single history—as a first-class concern for provenance, and argues
% that classical algebraic provenance implicitly assumes this ambiguity has already
% been resolved.

% \item It introduces \emph{interpretation provenance} as a semantic layer that
% records the commitments under which conclusions become well defined, and shows how
% classical algebraic provenance arises as a special case once interpretation is
% fixed.

% \item It formalizes \emph{layer certificates} as a minimal interface for carrying
% semantic commitments forward, enabling provenance reasoning that is sensitive to
% which distinctions a system chooses to make and which it intentionally collapses.

% \item It demonstrates, through two sandboxes, how interpretation provenance exposes
% previously implicit structure: in Datalog with negation, by making semantic
% dependency and resolution explicit; and in transactional systems, by explaining
% how query conclusions depend on ordering and isolation commitments.

% \item It identifies \emph{robustness} and \emph{semantic change} as new, well-posed
% classes of provenance questions—respectively, invariance across interpretations of
% a fixed specification and preservation across different specifications—and shows
% how interpretation provenance supports both.
% \end{enumerate}

% ================================================================
\section{Preliminaries: Ambiguous Semantics and Refinement}
\label{sec:prelim}
% ================================================================

Classical data provenance assumes a fixed, resolved semantics: for each input instance,
a specification (query) determines a single outcome, and provenance explains how that outcome was
derived.
In many systems, however, concurrency, negation, or underspecification admit
multiple admissible outcomes for the same input.

This section introduces a minimal semantic framework for reasoning about such
ambiguity.
We model specifications with multiple admissible outcomes and the semantic
commitments—called \emph{interpretations}—that restrict ambiguity until meaning
is determinate.
Our focus is purely structural: we introduce only the semantic objects needed to
define interpretation provenance and to reason about robustness and change across
interpretations, without assuming any particular execution model.

\subsection{Histories}

When semantics are ambiguous, it is no longer sufficient to model input as a
single static instance.
Instead, we require a representation of \emph{partial information}: what has
been observed, and which semantic distinctions remain unresolved.

We model such partial information using \emph{histories}.
A history records a collection of semantic events together with any known
precedence constraints among them.
Histories generalize familiar notions of input:
in logic programming, a history may consist of a set of ground facts;
in transactional systems, it may record transaction invocations, reads, writes,
and commits; in distributed settings, it may include message deliveries or
observations.
When no precedence constraints are present, histories collapse to unordered
sets and coincide with classical database instances.

\begin{definition}[History]
A \emph{history} is a finite partially ordered set
\[
  H = (E, \rightarrow),
\]
where $E$ is a set of events and $\rightarrow \subseteq E \times E$ is a strict
(irreflexive and transitive) partial order representing known precedence constraints among events.
\end{definition}

The order $\rightarrow$ records only irreducible constraints—facts that must be
respected by any admissible semantic interpretation.
Histories abstract semantic information from execution detail: they may coincide
with executions or traces, but they do not require a total order or a fixed
schedule, and apply equally to purely semantic settings such as Datalog with
negation.

We write $\Hist$ for the class of histories under consideration.
This choice is parametric and deliberately abstract: histories may represent
database updates, transaction invocations, message deliveries, or any other
events that constrain semantic interpretation.

\begin{definition}[History Extension]
For histories $H_1=(E_1,\rightarrow_1)$ and $H_2=(E_2,\rightarrow_2)$, we write
$H_1 \sqsubseteq H_2$ if:
\begin{enumerate}
  \item $E_1 \subseteq E_2$, and
  \item $\rightarrow_1 = \rightarrow_2 \cap (E_1 \times E_1)$.
\end{enumerate}
\end{definition}

Intuitively, histories are extended by accumulating events and precedence constraints,
without revising previously observed relationships.
We interpret $\sqsubseteq$ as information growth, not as time or execution \emph{per se}.

\subsection{Observations and Specifications}
Fix a set $O$ of fully specified semantic objects, called \emph{outcomes}.
\begin{definition}[Observation]
An \emph{observation} is an element of a set $O$ of fully specified semantic objects.
\end{definition}

Observations represent what an observer may conclude once all semantic commitments have
been fixed.
Examples include a database state, a serial execution of transactions, a model of
a logic program, the visible state of an automaton, or a completed task result.

\begin{definition}[Specification]
A \emph{specification} is a function
\[
  \Spec : \Hist \to 2^{O},
\]
mapping each history $H \in \Hist$ to a (possibly empty) set of admissible observations.
\end{definition}

The set $\Spec(H)$ captures all observations consistent with the evidence in $H$.
When $|\Spec(H)| > 1$, the specification admits multiple admissible meanings for $H$,
reflecting unresolved semantic ambiguity.
When $|\Spec(H)| = 1$, the meaning of $H$ is fully determined.

Ambiguity in this sense is common and intentional.
Specifications for Datalog with negation, transactional systems, and distributed
executions all admit multiple admissible observations for the same history,
corresponding to different semantic resolutions of underspecified behavior.
In such settings, semantic questions cannot be posed until ambiguity is constrained
by interpretation.\footnote{%
The ambiguity we study lives in the \emph{powerset domain} $2^{O}$: multiplicity
of admissible outcomes $\Spec(H) \subseteq O$.
This is distinct from uncertainty or nondeterminism \emph{within} the outcome
domain $O$ itself (e.g., distributions, multi-valued attributes, or
partial observations).
We discuss this distinction in more detail in Section~\ref{sec:negation} and
Appendix~\ref{sec:appendix-datalog}.}


\subsection{Semantic Refinement}

Ambiguity is constrained by making additional semantic commitments.
We model these commitments as \emph{interpretive predicates}, understood as
semantic conditions that \emph{filter} the set of admissible observations while
preserving feasibility.

\begin{definition}[Interpretive Predicate (as Outcome Filter)]
An \emph{interpretive predicate} is a function
\[
  \varphi : \Hist \times 2^{O} \to 2^{O}
\]
such that for all histories $H$ and observation sets $S \subseteq O$, \hspace{0.8mm} $\varphi(H,S)
\subseteq S$ (\emph{Shrinkage}), and if $S \neq \emptyset$ then $\varphi(H,S) \neq \emptyset$
 (\emph{Feasibility}).
\end{definition}
Interpretive predicates restrict admissible observations without introducing new ones or rendering
the specification infeasible.

Crucially, an interpretive predicate may depend on the admissible outcome set
\emph{as a whole}.
That is, $\varphi(H,S)$ need not be definable as a pointwise test applied
independently to each $o \in S$, but may instead enforce semantic commitments
whose effect depends on which alternatives are jointly admissible.
We refer to such predicates as \emph{ambiguity-sensitive}.

\begin{definition}[Refinement by Conditioning]
Given a specification $\Spec$ and an interpretive predicate $\varphi$, the
\emph{refinement} of $\Spec$ by $\varphi$ is the specification
\[
  (\Spec \mid \varphi)(H) \triangleq \varphi(H,\Spec(H)).
\]
\end{definition}

Multiple interpretive predicates may be applied by iterating this refinement
operation.
Because predicates may be ambiguity-sensitive, the order in which refinements
are applied can matter.

\subsection{Interpretations}

Individual interpretive predicates typically express only partial semantic
commitments.
We therefore model semantic resolution using finite compositions of refinement
predicates, interpreted as \emph{sequential conditioning}.

\begin{definition}[Interpretation]
An \emph{interpretation} is a finite expression
\[
  I = \varphi_1 \wedge \cdots \wedge \varphi_m
\]
of interpretive predicates, where $\wedge$ denotes \emph{ordered (sequential)
refinement} rather than logical conjunction.
In particular, $\wedge$ is not assumed to be commutative or idempotent.
\end{definition}

Refinement by an interpretation is defined by sequential application:
\[
  \Spec \mid I
  \;\triangleq\;
  (\cdots((\Spec \mid \varphi_1)\mid \varphi_2)\cdots)\mid \varphi_m .
\]

Because interpretive predicates may be ambiguity-sensitive, the order of conjuncts
can matter.
This non-commutativity induces semantic dependency.


\paragraph{Dependency graphs on interpretive predicates.}
A useful way to make this structure explicit is via a \emph{dependency graph}
over interpretive predicates.
An edge $\varphi \rightarrow \psi$ indicates that $\psi$ cannot be applied until
$\varphi$ has been resolved.
Such dependencies determine when a specification becomes resolved and when
depth--0 derivational provenance applies.

\begin{definition}[Resolved specification and resolving interpretation]
A specification $\Spec$ is \emph{resolved} if for every history $H$, $|\Spec(H)| = 1$.
An interpretation $I$ \emph{resolves} $\Spec$ (i.e., $I$ is \emph{resolving for} $\Spec$) if $\Spec \mid I$ is resolved.
\end{definition}

\paragraph{Clarification.}
Resolution is a property of the set-valued specification $\Spec:\Hist\to 2^O$:
it is \emph{single-valued pointwise in the history}, i.e., for every history $H$,
the set $\Spec(H)$ is a singleton.
Resolution marks the point at which semantic ambiguity has been fully discharged.

\subsection{Restriction and Identification}

Fix a specification $\Spec$ and an interpretation $I$.
For each history $H$, refinement by $I$ produces a restricted set of admissible
outcomes $(\Spec \mid I)(H)$.

The semantic effect of interpretation has two aspects:
\emph{restriction}, in which outcomes incompatible with the commitments in $I$ are eliminated and play no further semantic role, and \emph{identification}, in which some of the remaining admissible outcomes become semantically indistinguishable.

Formally, for a history $H$ write
\[
  O_H^I \triangleq (\Spec \mid I)(H).
\]
We say that two outcomes $o,o' \in O_H^I$ are \emph{equivalent under $I$ at $H$},
written $o \sim_{I,H} o'$, if no interpretive predicate in $I$ distinguishes between them at $H$.

\begin{definition}[Interpretation-Induced Quotient]
For each history $H$, interpretation $I$ induces a quotient
\[
  O_H^I \big/ {\sim_{I,H}},
\]
obtained by identifying admissible outcomes that are equivalent under $I$.
\end{definition}

The equivalence relation $\sim_{I,H}$ is total on the restricted domain $O_H^I$.
Outcomes ruled out by the interpretation are not in $O_H^I$, and hence simply absent from the quotient.

\begin{remark}[History-indexed quotients]
Because interpretive predicates may depend on the history and on intermediate
admissible sets, equivalence under an interpretation is defined relative to a
history.
Interpretations therefore induce a family of quotients indexed by histories,
rather than a single global quotient on $O$.
\end{remark}

These notions of refinement, restriction, and identification provide the raw
semantic substrate on which we build algebraic structure in the next section.

%% ================================================================
\section{Interpretations and Provenance Structure}
\label{sec:conditional-semiring}
This section connects the semantic notion of interpretation
(Section~\ref{sec:prelim}) to provenance.
Fix a specification $\Spec$ and a resolving interpretation $I$.
Interpretation provenance records both the semantic commitments fixed by $I$ and
the derivational provenance that applies after conditioning on $I$.

\subsection{Syntax of Interpretation Provenance}

We now introduce a compact notation for \emph{interpretation provenance}.
This subsection introduces no new semantic machinery.
It packages together notions already defined—interpretations, resolution, and
classical provenance—into a single explanatory object.
Throughout the paper, interpretation provenance is expressed by conditioning
semantic objects on an interpretation.
At this stage, we assume no algebraic structure on interpretations or outcomes
beyond what has already been defined.

\paragraph{Interpretation provenance values.}
Fix a commutative semiring $(K,\otimes,1)$ for derivational provenance
(e.g., polynomial provenance).
For a query $Q$ and tuple $t$, write $\Prov_{Q \mid I}(t) \in K$ for the
ordinary (depth--0) provenance computed under the conditioned specification
$\Spec \mid I$.
We define \emph{interpretation provenance} as the pair
\[
  \mathsf{IP}_Q(t) \;\triangleq\; (I,\, \Prov_{Q \mid I}(t)).
\]

% When convenient, we may write $I \cdot k$ as an abbreviation for the pair $(I,k)$,
% but all semantics are expressed via conditioning on $I$.

The intended reading is:
\emph{$k$ explains the derivation, assuming the semantic commitments fixed by $I$.}

\paragraph{Relationship to classical provenance.}
Once an interpretation $I$ resolves $\Spec$, semantic ambiguity is fixed for the purpose of provenance.
Query evaluation under $\Spec \mid I$ is single-valued, and classical semiring
provenance applies without modification.
Interpretation provenance therefore records classical provenance under
conditioning, together with the interpretation that justifies that conditioning.

In particular, for \emph{depth--0} specifications—where no nontrivial semantic
commitments are required—interpretation provenance coincides exactly with
classical semiring-based provenance.

\paragraph{Why interpretation provenance is interpretation-relative.}
Interpretation provenance is defined relative to a \emph{single, fixed}
interpretation.
A provenance value records how a semantic outcome is derived \emph{given} a
particular set of semantic commitments; it does not combine or quantify over
multiple interpretations.

Accordingly, we do not form disjunctions or aggregations over interpretations
within a single provenance object.
Questions about how outcomes or explanations vary across different
interpretations—such as robustness, invariance, or sensitivity to semantic
change—are treated explicitly and separately in
Section~\ref{sec:robustness}.

We illustrate this syntax next in a transactional setting, where
different ordering commitments induce different interpretations and,
consequently, different provenance for the same query.

\subsection{Example Revisited: Resolving Transactional Ambiguity}
\label{sec:concurrency}

We now return to the transactional Example~\ref{ex:example}.
Recall that values $b, c$ are present in every serializable execution, while value $d$
is present only under certain execution orders.

We can now show the syntax of interpretation provenance for this example.
In a conflict serializability context, interpretive predicates correspond to ordering
commitments among transactions: $\varphi_{T_i \prec T_j}$ commits to
ordering the actions of transaction $T_i$ before the conflicting actions of $T_j$.
If we observe $\{(b), (c), (d)\}$ as the query result, then an example Interpretation is 
\[
I_{\mathsf{in}} \;\triangleq\; \varphi_{T_1 \prec T_2} \wedge \varphi_{T_2 \prec T_3} \wedge \varphi_{T_3 \prec T_4} \wedge \varphi_{T_4 \prec T_Q} \wedge \varphi_{T_Q \prec T_5}
\]
and the corresponding conditioned (depth--0) provenance is
\[
  \Prov_{Q \mid I_{\mathsf{in}}}(b)=x_1+x_2,\;
  \Prov_{Q \mid I_{\mathsf{in}}}(c)=x_3,\;
  \Prov_{Q \mid I_{\mathsf{in}}}(d)=x_4.
\]

Interpretation provenance pairs these expressions with the interpretation,
e.g., $\mathsf{IP}_Q(b) \triangleq (I_{\mathsf{in}},\, x_1+x_2)$.

In Appendix~\ref{sec:appendix-concurrency} we work this example more fully, including 
the case where $d$ is absent from the query result.

\subsection{Interpretation Parsimony and Discriminating Power}
\label{sec:parsimony}

Recall that for a history $H$ and interpretation $I$, refinement induces a
quotient $O_H^I \big/ {\sim_{I,H}}$,
which records which admissible outcomes remain semantically distinguished after
conditioning (Section~\ref{sec:prelim}).

\paragraph{Discriminating power.}
Fix a specification $\Spec$.
An interpretation $I$ \emph{discriminates} between two outcomes
$o,o' \in \Spec(H)$ if $o \not\sim_{I,H} o'$.
If $o \sim_{I,H} o'$ for all resolving interpretations, then the available
semantic commitments never separate those outcomes.

\paragraph{Parsimonious predicate bases.}
Let $\mathcal{B}$ be a basis of interpretive predicates.
We call $\mathcal{B}$ \emph{parsimonious} if:
(i) for every history $H$, some interpretation formed from $\mathcal{B}$
resolves $\Spec$ at $H$; and
(ii) for every $H$ and every pair of distinct admissible outcomes
$o \neq o' \in \Spec(H)$, there exists a resolving interpretation from
$\mathcal{B}$ that discriminates between them.

Intuitively, a parsimonious basis induces quotients that are no coarser than
necessary: any semantic distinction that exists in principle can be exposed by
some admissible interpretation.

\paragraph{Parsimony and robustness.}
Parsimony is a property of the \emph{predicate basis}, not of individual queries.
The predicate basis determines which semantic distinctions are expressible, and
therefore which admissible outcomes can be distinguished or collapsed.
In practice, systems often adopt non-parsimonious bases—for reasons of efficiency,
policy, or usability—thereby intentionally abstracting away fine-grained
differences among admissible outcomes.

In the transactional setting, for example, predicates that commit only to an
equivalent serial schedule are deliberately coarse.
They collapse many concrete execution orders into a single semantic class, leaving
some tuples—such as $a$ and $b$ in Section~\ref{sec:intro}—present under every
interpretation consistent with that abstraction.
Interpretation provenance makes this explicit by showing which observations are
invariant across all interpretations induced by the chosen predicate basis, and
which would require sharper ordering predicates to distinguish.

This invariance is not accidental.
Observations are robust precisely when the available interpretive vocabulary
never discriminates between the outcomes on which they would differ.
This perspective motivates the robustness questions studied in
Section~\ref{sec:robustness}, where robustness is understood as persistence
\emph{relative to a fixed predicate basis}, rather than as an absolute semantic
property.

\subsection{Summary}
Interpretation provenance pairs a resolving interpretation with classical
provenance computed under that interpretation.
Once an interpretation is fixed, provenance is classical by construction: it is
computed over the conditioned specification $\Spec \mid I$.

The following sections instantiate this framework in two settings with very
different sources of semantic ambiguity: negation in Datalog and scheduling in
transactional systems.


% ================================================================
\section{Layer Certificates and Interpretation Provenance}
\label{sec:certificates}
% ================================================================

Interpretations resolve semantic ambiguity by conditioning a specification.
In this section we ask a complementary question:

\begin{quote}
\emph{What information must be retained to support further interpretation
\emph{and} provenance queries about the result?}
\end{quote}

Interpretation proceeds incrementally, but neither subsequent interpretation nor
provenance analysis requires access to the full refined specification.
Instead, it suffices to retain summaries that preserve exactly the semantic
distinctions relevant to future reasoning.
We formalize these summaries as \emph{layer certificates}.

\subsection{Certificates and Provenance Queries}

Fix a specification $\Spec$ and an interpretation
$I = \varphi_1 \wedge \cdots \wedge \varphi_k$ with layering
$L_1,\ldots,L_k$.
Let
\[
  \Spec_i \;\triangleq\; \Spec \mid (\varphi_1 \wedge \cdots \wedge \varphi_i),
  \qquad
  J \;\triangleq\; \varphi_{i+1} \wedge \cdots \wedge \varphi_k .
\]

\paragraph{Certificates.}
An interpretive predicate restricts admissible outcomes.
A \emph{certificate} summarizes this restriction.

\begin{definition}[Certificate]
Let $\varphi$ be an interpretive predicate.
A value $C$ with denotation $\llbracket C \rrbracket \subseteq O$ is a
\emph{certificate} for $\varphi$, written $\varphi \models C$, if for all histories
$H$ and nonempty $S \subseteq O$,
\[
  \varphi(H,S) \subseteq S \cap \llbracket C \rrbracket .
\]
\end{definition}

Certificates may be weaker than the predicates they summarize: they need only
exclude outcomes that would be excluded by interpretation, not encode the full
resolution process.

\paragraph{Provenance queries.}
Interpretation provenance is queried.

\begin{definition}[Provenance Query]
A \emph{provenance query} is a function
\[
  q : (\Spec \mid I) \to A ,
\]
where $A$ is an answer domain (e.g., Boolean, explanations, or dependency sets).
\end{definition}

Examples include robustness queries (“does $t$ hold under all interpretations in a
class?”), inter-specification semantic-change queries (Section~\ref{sec:robustness}), and
conditional why/why-not queries.

\paragraph{Sufficiency for continuation and queries.}
Certificates must preserve both future interpretation and the answers to intended
provenance queries.

\begin{definition}[Sufficient Certificate]
A certificate $C_i$ is \emph{sufficient for $(J,\mathcal{Q})$} if for all histories
$H$:
\begin{enumerate}
  \item \emph{Continuation:}
  \[
    (\Spec_i \mid J)(H)
    =
    \bigl((\Spec \mid \chi_{C_i}) \mid J\bigr)(H);
  \]
  \item \emph{Query preservation:} for all $q \in \mathcal{Q}$,
  \[
    q\bigl((\Spec_i \mid J)(H)\bigr)
    =
    q\bigl(((\Spec \mid \chi_{C_i}) \mid J)(H)\bigr).
  \]
\end{enumerate}
\end{definition}

A \emph{layer certificate} is a certificate sufficient for all subsequent layers
and the intended family of provenance queries.

\subsection{Query Sensitivity of Certificates}

Supporting richer provenance queries strictly increases certificate requirements.

\begin{theorem}[Query Monotonicity]
If $\mathcal{Q}_1 \subseteq \mathcal{Q}_2$, then every certificate sufficient for
$(J,\mathcal{Q}_2)$ is sufficient for $(J,\mathcal{Q}_1)$.
\end{theorem}

\noindent
Thus, extending the space of provenance questions can only require retaining
additional semantic information, never less.

Concrete illustrations of this phenomenon appear in the appendices.
Appendix~\ref{sec:appendix-datalog} shows how inter-specification semantic-change queries
(Section~\ref{sec:robustness}) for negation semantics enlarge certificates beyond
those needed for derivational provenance.
Appendix~\ref{sec:appendix-concurrency} shows the analogous effect for robustness
queries under transactional ambiguity.

\subsection{Summary}

Layer certificates provide a principled interface between semantic resolution and
provenance.
By parameterizing certificates by families of provenance queries,
interpretation provenance makes explicit which semantic distinctions must be
retained—and which can be discarded—when supporting explanation, robustness, or
inter-specification semantic change (Section~\ref{sec:robustness}).

This perspective also clarifies a familiar pattern in systems lineage and
observability.
Tracing systems such as Dapper typically retain worst-case information in order to
support arbitrary downstream diagnostic queries, effectively treating all
semantic distinctions as potentially relevant~\cite{sigelman2010dapper}.
In our framework, this corresponds to choosing certificate families large enough
to preserve all future provenance queries.
Interpretation provenance exposes this tradeoff explicitly and provides a semantic
basis for more parsimonious lineage when the intended query family is known.
We illustrate this connection in Appendix~\ref{app:systems-lineage}.

% ================================================================
\section{Sandbox I: Datalog with Negation As Semantic Ambiguity}
\label{sec:negation}
% ================================================================

We instantiate interpretation provenance in a familiar and deeply studied
setting: Datalog with negation.
Our purpose is not to propose new semantics for negation, nor to compare their
expressive power.
Rather, we use classical semantics for Datalog with negation as a controlled
environment in which to observe how semantic ambiguity in a \emph{fixed}
specification gives rise to different structures of interpretation provenance.

Datalog with negation exhibits ambiguity even for a fixed program and a fixed
input instance.
Negated predicates introduce dependencies that delay semantic commitment: whether
an atom is justified can depend on the absence of other atoms whose status is not
yet determined.
Interpretation provenance makes these dependencies explicit, recording not only
what conclusions are drawn, but the semantic commitments under which they are
valid.

We examine three standard semantics for negation—stratified, well-founded, and
stable—and write down the corresponding interpretation provenance in each case.
Figure~\ref{fig:interp-tables}(A) summarizes these regimes as interpretation tables,
making explicit which interpretive predicates are sufficient to resolve the
specification under each semantics.


\begin{figure*}[t]
\centering
\footnotesize
\setlength{\tabcolsep}{6pt}
\renewcommand{\arraystretch}{1.15}

\begin{tabular}{
@{}p{0.12\linewidth} p{0.30\linewidth} || 
   p{0.24\linewidth} p{0.26\linewidth}@{}
}
\toprule
\multicolumn{2}{c||}{\textbf{(A) Datalog negation semantics}} &
\multicolumn{2}{c}{\textbf{(B) Transactions under serializability}} \\
\midrule
\textbf{Case} & \textbf{Resolving interpretation(s)} &
\textbf{Case} & \textbf{Minimal interp. predicates} \\
\midrule
Stratified &
$\varphi_{\mathrm{stab}(p)} \wedge \varphi_{\mathrm{neg}(q)}$ &
SER robust $b$ &
$\varphi_{T_1 \prec T_Q} \wedge \varphi_{T_2 \prec T_Q}$ \\
%
Well-founded &
$\varphi_{p(c)=\mathbf{t}} \wedge
 \varphi_{q(c)=\mathbf{f}} \wedge
 \varphi_{r(c)=\mathbf{u}} \wedge
 \varphi_{s(c)=\mathbf{u}}$ &
SER contingent $d$ (why) &
$\varphi_{T_4 \prec T_Q} \wedge \varphi_{T_Q \prec T_5}$ \\
%
Stable &
$I_{\mathrm{stab}}^{(r)} \;\text{or}\;
I_{\mathrm{stab}}^{(s)}$ (incompatible) &
SER contingent $d$ (why-not) &
$\varphi_{T_5 \prec T_Q}$ \\
\bottomrule
\end{tabular}

\caption{
Interpretation summaries for our two sandboxes.
(\textbf{A}) In Datalog, negation semantics differ in whether they admit a single
resolving interpretation (stratified, well-founded) or multiple incompatible ones
(stable).
(\textbf{B}) Under serializability, many serial orders are equivalent; we therefore
list only the minimal interpretive predicates required to justify a query result. 
Full details are given in Appendices~\ref{sec:appendix-concurrency} and ~\ref{sec:appendix-datalog}.
}
\Description{
A single side-by-side table comparing interpretation structure.
Left: resolving interpretations for Datalog negation semantics.
Right: minimal interpretive predicates needed to justify transactional query results
under serializability.
}
\label{fig:interp-tables}
\end{figure*}

Full worked interpretations, dependency graphs, and conditioned provenance values
appear in Appendix~\ref{sec:appendix-datalog} and Appendix~\ref{sec:appendix-concurrency}.
These distinctions will reappear in the next section, where ambiguity arises from
transaction schedules rather than negation.

\subsection{Specification and Semantic Ambiguity}

A Datalog program with negation induces a specification
$\Spec : \Hist \to 2^{O}$ mapping histories (finite sets of ground EDB atoms over the fixed constant $c$ introduced below) to sets of admissible
outcomes (models).
Ambiguity arises because negation can create circular semantic dependencies:
whether an atom holds may depend on whether some other atom does not hold, and
vice versa.

\paragraph{Running example.}
Consider the following program $P$ over EDB facts in a history $H$.
For notational economy we use a fixed constant $c$ and take all predicates unary, so the EDB fact is $a(c)$ and the derived atoms are $p(c),q(c),r(c),s(c)$.
\begin{align*}
  p(c) &\leftarrow a(c).\\
  q(c) &\leftarrow \neg p(c).\\[1mm]
  r(c) &\leftarrow \neg s(c).\\
  s(c) &\leftarrow \neg r(c).
\end{align*}
Assume $H$ contains the single extensional fact $a(c)$.
The pair $(P,H)$ contains both a stratified fragment
($q(c)$ depends negatively on $p(c)$, and $p(c)$ depends positively on $a(c)$)
and a cycle through negation ($r \leftrightarrow \neg s$).
We use this single instance to exhibit, side-by-side, three distinct structures
of interpretation provenance.

\subsection{Interpretive Predicates for Negation}

Across all negation semantics we consider, interpretation proceeds by applying
interpretive predicates that constrain admissible models.
In this sandbox, it is helpful to think of interpretive predicates as semantic
commitments about the status of atoms---e.g., committing that an atom is true,
false, or (when the outcome domain is lifted) undefined.

These commitments are not local to individual derivations.
They depend on which outcomes remain admissible and may enable or disable further
commitments.
As a result, interpretive predicates for negation are ambiguity-sensitive and may
fail to commute.

\subsection{Stratified Semantics}

Stratified semantics applies to programs whose dependency graph contains no cycles
through negation.
It restricts attention to specifications for which negation can be interpreted by
a stratum order, allowing semantic commitments to be staged.

In our running example, we temporarily set aside the predicates $r(c)$ and $s(c)$,
which form a negation cycle, and focus on the stratified fragment defining $p(c)$
and $q(c)$.
In this fragment, $p(c)$ lies in a lower stratum derived directly from the EDB,
while $q(c)$ lies in a higher stratum, depending negatively on $p(c)$.

The central semantic issue is not the computation of a fixpoint per se, but the
commitment that the lower stratum has \emph{stabilized}—that is, that no further
lower-stratum consequences will appear that could change the truth of $p(c)$.
In incremental or distributed settings, this stabilization is a genuine semantic
choice: the possibility of late-arriving facts creates ambiguity about whether
negation is safe to apply, and resolving that ambiguity requires committing to a
closed world for the lower stratum.

Accordingly, resolution proceeds by two staged interpretive predicates:
$\varphi_{\mathrm{stab}(p)}$, committing that the lower stratum defining $p(c)$ has
stabilized, followed by $\varphi_{\mathrm{neg}(q)}$, committing to the stratified
interpretation of $\neg p(c)$.
The dependency $\varphi_{\mathrm{stab}(p)} \rightarrow \varphi_{\mathrm{neg}(q)}$
yields a single resolving interpretation, as summarized in the first row of
Figure~\ref{fig:interp-tables}(A).
The fully resolved specification supports ordinary derivational provenance for
$q(c)$ under this interpretation.
The explicit interpretation $I_{\mathrm{strat}}$, along with its dependency
structure and conditioned provenance values, is worked out in
Appendix~\ref{sec:appendix-datalog}.

\subsection{Well-Founded Semantics}

Well-founded semantics (WFS) applies to all Datalog programs with negation,
including those with cycles.
Rather than selecting among multiple total two-valued models, WFS resolves
ambiguity by enlarging the outcome domain: atoms may be \emph{true}, \emph{false},
or \emph{undefined}.
Undefinedness is not a failure of semantics; it is an explicit representation of
residual ambiguity \emph{within the outcome domain}.
Rather than resolving multiplicity among admissible outcomes, well-founded
semantics internalizes ambiguity by lifting the domain of outcomes itself.
This highlights a modeling distinction that our framework makes explicit.
Ambiguity in $\Spec(H) \subseteq O$ corresponds to multiple incompatible semantic
outcomes and is resolved by interpretation.
By contrast, lifting the outcome domain (as in well-founded semantics) moves
uncertainty \emph{into} $O$ itself, yielding outcomes that are intentionally
multi-valued and meaningful to the application rather than requiring further
interpretation.

For our example, the stratified prefix remains unchanged: the predicates
$\varphi_{\mathrm{stab}(p)}$ and $\varphi_{\mathrm{neg}(q)}$ are applied as before.
The remaining cycle between $r(c)$ and $s(c)$ induces a mutual dependence that
cannot be broken by staging alone.
Under well-founded semantics, this dependence is resolved deterministically by
committing both atoms to the value $\mathbf{u}$.

The resulting interpretation is unique as an element of $2^O$: there is a single
admissible outcome.
Resolution is achieved not by choosing among alternatives, but by lifting the
outcome domain to include an explicit undefined value.
This structure is reflected in the second row of
Figure~\ref{fig:interp-tables}(A).
Appendix~\ref{sec:appendix-datalog} presents the explicit resolving interpretation
$I_{\mathrm{wfs}}$, its dependency structure, and the resulting conditioned and
interpretation provenance values.

\subsection{Stable Model Semantics}

Stable model semantics resolves negation cycles in a fundamentally different way.
Rather than admitting partial outcomes, it insists on total two-valued models and
therefore requires an explicit semantic choice when cycles arise.

In our example, the staged prefix again applies:
$\varphi_{\mathrm{stab}(p)} \rightarrow \varphi_{\mathrm{neg}(q)}$.
However, the cycle between $r(c)$ and $s(c)$ now admits multiple incompatible
resolutions.
Each stable model corresponds to a distinct resolving interpretation, committing
one atom to true and the other to false.
No single interpretation refines the others; the admissible outcomes are separated
into mutually exclusive semantic worlds.

This multiplicity of resolving interpretations is summarized in the third row of
Figure~\ref{fig:interp-tables}(A).
Interpretation provenance must therefore record \emph{which} stable interpretation
was assumed before derivational provenance can be applied.
Appendix~\ref{sec:appendix-datalog} gives the explicit interpretations, dependency
graphs, and provenance values for each stable model.

\subsection{Takeaways for Interpretation Provenance}

This sandbox highlights three distinct provenance shapes induced by three
classical negation semantics.
Stratification yields a layered structure in which negation-dependent derivations
are conditional on stabilization commitments for earlier strata.
Well-founded semantics resolves ambiguity by lifting the outcome domain and
recording explicit commitments to partiality (undefinedness) in the presence of
negation cycles.
Stable semantics yields multiple resolving interpretations for the same program
and instance, so interpretation provenance must record which resolved
interpretation was assumed even before derivational provenance is consulted.
When negation is acyclic, all three semantics coincide and interpretation
provenance collapses to ordinary derivational provenance, reflecting the absence
of semantic ambiguity.

Appendix~\ref{sec:appendix-datalog} gives the explicit interpretations, dependency graphs, and provenance values.

\paragraph{Semantic change (forward reference).}
Semantic change is an inter-specification question; see Section~\ref{sec:robustness}.
A worked semantic-change analysis appears in Appendix~\ref{sec:appendix-datalog}.

The next sandbox mirrors these interpretation shapes, but with ambiguity arising
from concurrency rather than negation.

% ================================================================
\section{Sandbox II: Transactions as Systems Ambiguity}
\label{sec:transactions}
% ================================================================

We now turn to transactional systems.
As in the previous section, our goal is not to propose new correctness criteria,
but to understand how semantic ambiguity shapes provenance.

Transactional systems introduce a qualitatively different source of ambiguity
from negation.
Here, ambiguity arises from \emph{concurrency}: transactions race to read and write
shared state, and the execution may not uniquely determine which effects are
observed.
Interpretation provenance explains how query results depend not only on data, but
on the semantic resolution of that concurrency.

\subsection{Provenance Questions Under Concurrency}

Classical provenance explains how an output tuple depends on base data, assuming a
fixed semantics and a fixed execution outcome.
Transactional systems invalidate this assumption.

Consider a query result tuple $t$ produced by executing a workload of
transactions.
There are several distinct provenance questions one might ask: (\emph{Existence})
does $t$ appear under \emph{some} admissible execution; (\emph{Universality}) does
$t$ appear under \emph{all}
admissible executions; and (\emph{Dependence}) on which semantic choices---e.g.,
transaction orderings or visibility decisions---does the presence of $t$ depend?

These questions are strictly more subtle than asking whether a final database
state is invariant across executions.
A database state may vary across admissible schedules even when a particular
query answer is stable; conversely, a state may be fixed while a query result is
not.

Answering these questions requires provenance that is sensitive to
\emph{which transactional resolutions were relied upon}, not merely which data
items were read.
This is precisely the role of interpretation provenance.

\subsection{Ambiguity Before Order}

To address this setting, we begin with a deliberately
weak semantic model.
A transactional workload consists of a set of transactions, each issuing reads
and writes on shared objects and eventually committing.
At this level, the execution does not specify an order between transactions.
Reads may observe values written by different transactions, and multiple
explanations of the same observed behavior may coexist.

This lack of order is the source of semantic ambiguity.
Even with complete knowledge of all operations, multiple outcomes may be
admissible.
Order is not a primitive feature of the execution; it is a semantic commitment
used to justify an outcome.

Classical provenance assumes a fixed execution outcome.
Transactional systems invalidate this assumption: the same history may admit
multiple outcomes, and tuple-level provenance questions must therefore range over
semantic resolutions of concurrency.

\subsection{Outcomes under Serializability}

Under serializability, outcomes are defined via serialization.

An \emph{outcome} is a final database state produced by some serial execution of
the transactions.
An outcome is admissible if there exists a total order of transactions such that
sequential execution in that order yields the same final state and respects the
history’s ordering constraints.

A single history may admit many such outcomes.
Each corresponds to a distinct semantic explanation of the same concurrent
behavior.
Interpretation constrains this space by committing to ordering decisions.

\subsection{Interpretive Predicates as Ordering Commitments}

Accordingly, the primitive interpretive predicates in this setting are
\emph{ordering predicates}.
For transactions $T_i$ and $T_j$, the predicate
$\varphi_{T_i \prec T_j}$ commits that $T_i$ precedes $T_j$ in the serialization.

Conditioning on an ordering predicate eliminates all admissible outcomes whose
serializations violate it.
An interpretation is a finite conjunction of such predicates, inducing a partial
order on transactions.
A fully resolving interpretation induces a total order; partial interpretations
leave multiple serializations admissible.

As in the negation setting, ordering predicates are ambiguity-sensitive.
Their effect depends on which serializations remain admissible, and different
ordering commitments may enable or preclude others.
This non-commutativity induces semantic dependency and interpretation depth.

\subsection{Interpretation Structure (Running Example)}

We reuse the running example from the introduction and
Appendix~\ref{sec:appendix-concurrency}, focusing on the contingent tuple $d$
inserted by $T_4$ and deleted by $T_5$.

Under the interpretation
\[
  I_{\mathsf{in}} \;\triangleq\;
  \varphi_{T_4 \prec T_Q} \wedge \varphi_{T_Q \prec T_5},
\]
the insert precedes the query and the delete follows it.
Conditioning on $I_{\mathsf{in}}$ yields a resolved snapshot in which $d$ is
present, and classical provenance applies:
\[
  \Prov_{Q \mid I_{\mathsf{in}}}(d)=x_4.
\]

Under the alternative interpretation
\[
  I_{\mathsf{out}} \;\triangleq\; \varphi_{T_5 \prec T_Q},
\]
the delete precedes the query, so $d$ is absent and the appropriate classical
question is \emph{why-not}.

Prior to fixing an interpretation, there is no single classical provenance
question for $d$.
Interpretation provenance makes this dependence explicit by recording the semantic
commitments under which derivations hold:
\[
  \mathsf{IP}_Q(d) \;\triangleq\; (I_{\mathsf{in}},\, x_4).
\]

Thus, classical provenance explains \emph{why} a tuple follows from a resolved
outcome, while interpretation provenance explains \emph{under which semantic
resolutions} that explanation applies.

\subsection{Beyond Serializability}

Other isolation levels expose the same phenomenon with different interpretive
predicates.
Snapshot isolation, for example, restricts admissible visibility relations while
preserving ambiguity; weaker isolation levels admit even larger families of
outcomes.

Across all regimes, the provenance question is unchanged:
which semantic commitments were required to justify a result?
Interpretation provenance provides a uniform account.
Mixed-regime systems (e.g., SER combined with SI) are treated as an explicit
policy choice in Section~\ref{sec:robustness}.

\paragraph{Summary.}
Transactional ambiguity arises because order is not given but must be assumed.
Interpretation provenance isolates the semantic commitments that justify query
results, extending classical provenance to settings where correctness admits
multiple valid executions.

% ================================================================
\section{Change and Robustness Across Interpretations}
\label{sec:robustness}
% ================================================================

Interpretation provenance records not only how observations are derived, but the
semantic commitments under which those derivations are valid.
This section formalizes two such families of questions and grounds them in the
examples developed earlier.

Here, \emph{observations} refer to observable query answers or derived facts—i.e., elements of the outcome set $O$—as defined in Section~\ref{sec:prelim}.

\subsection{Intra- vs.\ Inter-Specification Questions}

We distinguish two classes of provenance questions.

\emph{Intra-specification} questions fix a specification $\Spec$ and range over its
resolving interpretations.
They ask which observations hold under different admissible resolutions
of the \emph{same} specification.

\emph{Inter-specification} questions compare observations across distinct
specifications $\Spec$ and $\Spec'$.
They ask which observations persist when the specifications themselves change---for
example, when a transactional system moves from serializability to snapshot
isolation, or when a logic program is interpreted under a different negation
semantics.

Both question families are expressed using interpretation provenance, they
just quantify over different spaces.

\subsection{Robustness Within a Fixed Specification}

Robustness is an intra-specification property.
Fix a specification $\Spec$ and a class $\mathcal{I}$ of resolving interpretations.
An observation $o$ is \emph{robust} with respect to $\mathcal{I}$ iff
\[
  \forall I \in \mathcal{I}.\; o \in (\Spec \mid I)(H),
\]
where (recall Section~\ref{sec:prelim}) each $I \in \mathcal{I}$ is resolving, i.e., $\Spec \mid I$ is single-valued for each history.

Equivalently, $o$ does not depend on semantic commitments that vary across
$\mathcal{I}$.
In the lens of Section~\ref{sec:parsimony}, robustness is evidence of a lack of
parsimony in the interpretive predicate basis.

\paragraph{Example: robustness under serializability.}
In the transactional example of
Appendix~\ref{sec:appendix-concurrency}, the history admits multiple serial orders.
Interpretive predicates record only ordering commitments between transactions.
Under all such orders, the query result includes the values $b$ and $c$, while the
presence of $d$ depends on whether the insert $T_4$ is ordered before the query
and the delete $T_5$ after it.

Our interpretive predicate basis captures the level of discrimination natural to
serializability: it distinguishes executions by transaction order, but not by
finer-grained, value-specific effects.
This exposes that $b$ and $c$ are robust under serializability, while $d$ is not.
The lack of parsimony in the predicate basis here is both meaningful (it exposes artifacts of serializability)
and practical (it avoids an explosion of ordering predicates).

\subsection{Semantic change across specifications}
\label{sec:semantic-change}

Semantic change is an inter-specification question: it compares observations
admissible under different specifications $\Spec$ and $\Spec'$.

\paragraph{Preservation.}
Given specifications $\Spec$ and $\Spec'$, an outcome $o$ is \emph{preserved}
from $\Spec$ to $\Spec'$ iff
\[
  \forall I \in \mathcal{I},\; \forall I' \in \mathcal{I}'.\;
  \bigl(o \in (\Spec \mid I)(H)\bigr)
  \;\Rightarrow\;
  \bigl(o \in (\Spec' \mid I')(H)\bigr),
\]
where $\mathcal{I}$ and $\mathcal{I}'$ range over resolving interpretations
admissible under $\Spec$ and $\Spec'$.

\paragraph{Example: serializability versus snapshot isolation.}
Extend the running example on relation $S(k,y)$ with two additional transactions:
\[
\begin{array}{ll}
T_6: \textsf{read } S(5,e);\ \textsf{delete } S(6,f) \qquad &
T_7: \textsf{read } S(6,f);\ \textsf{delete } S(5,e).
\end{array}
\]
Assume both tuples $S(5,e)$ and $S(6,f)$ are initially present.

Under snapshot isolation, both transactions may read the same snapshot and both
deletes may commit.
The outcome in which neither $e$ nor $f$ appears in the query result is therefore
admissible.
That outcome is not admissible under serializability and is therefore
\emph{not preserved} when strengthening the specification from snapshot isolation
to serializability.

By contrast, every serializable execution is admissible under snapshot isolation.
Thus all outcomes admitted under serializability—including the robust query
results $b$ and $c$ from the original workload—are preserved under the weakening
from serializability to snapshot isolation.
The tuple $d$ remains interpretation-dependent under both specifications.

\paragraph{Mixed-spec execution as policy.}
Many database systems allow isolation levels to be changed over time.
If isolation is ``tightened'' (e.g., from SI to SER), preserved outcomes remain valid, but
non-preserved outcomes (such as the joint absence of $e$ and $f$) must be dealt with according
to some policy---e.g. mask them as missing, reject transactions that touch them, or
carry interpretation provenance annotations through all queries so applications can decide how
to handle the ``taint''—i.e., the lack of preservation—of non-preserved results.
Interpretation provenance supports any such policy, by making the semantic commitments
behind each outcome explicit and checkable.

\subsection{Quantitative Variants}

The definitions above are qualitative.
When the space of interpretations is finite or otherwise measurable, one can ask
quantitative questions—for example, what fraction of admissible interpretations
justify a conclusion, or how often a conclusion is preserved across semantic
changes.
Interpretation provenance provides the semantic structure needed to formulate such
questions, though it does not require a measure over interpretations.

\paragraph{Summary.}
Robustness identifies conclusions invariant under the semantic distinctions a
system chooses to make.
Semantic change compares conclusions across different semantic regimes.
Interpretation provenance provides a common framework for expressing both, while
making explicit the limits of discrimination imposed by the chosen interpretive
predicates.

% ================================================================
\section{Related Work}
\label{sec:related}
% ================================================================

\subsection{Data Provenance and Explanations}

Algebraic provenance provides a uniform framework for explaining query
conclusions under fixed, resolved semantics by annotating derivations with
elements of a commutative semiring
\cite{green2007provenance,suciu2011probabilistic,dalvi2012probabilistic,olteanu2015factorized,senellart2018provenance}.
A related body of work studies explanations and causality for query answers,
including why-not provenance, causal responsibility, and counterfactual
explanations
\cite{buneman2001why,halpern2001causes,meliou2010causality,glavic2015survey}.

From the perspective of this paper, these approaches operate at
interpretation depth~0: semantic ambiguity has already been resolved, and
provenance explains derivations relative to a fixed semantics.
Interpretation provenance is complementary.
It explains which semantic commitments are required for observations to hold
at all, and how classical provenance depends on those commitments once
resolution has occurred.

\subsection{Negation and Non-Monotonic Semantics}

Classical semantics for Datalog with negation provide well-understood mechanisms
for resolving semantic ambiguity, with established relationships under syntactic
restrictions.
We do not propose new semantics.
Instead, Section~\ref{sec:negation} uses Datalog with negation as a controlled
setting in which resolution strategies are explicit and formal, making it
possible to study how different semantic commitments induce distinct
interpretation structures and provenance behavior.

\subsection{Transactional and Systems Provenance}

A first line of work studies provenance in transactional, workflow, and database
settings under an assumed execution or isolation level
\cite{bernstein1987concurrency,adya1999weak,cerone2015consistency,cheney2009provenance,freire2008provenance}.
In these settings, semantic resolution is implicit: provenance explains results
relative to a chosen execution, snapshot, or schedule.
From our perspective, such approaches again operate at interpretation depth~0.

A distinct body of systems work studies execution provenance or systems lineage,
recording causal dependencies among processes, messages, and events to explain
observed behavior
\cite{seltzer2005provenance,bates2015trustworthy,sigelman2010dapper}.
These approaches take a maximalist view of provenance, capturing entire execution
histories to account for ambiguity in scheduling and message interleaving.

Interpretation provenance provides a semantic foundation for structuring such
systems provenance.
Rather than committing to a single reconstructed execution, it makes explicit
the semantic commitments required for a diagnostic conclusion to hold across
admissible executions, supporting principled notions of parsimony and robustness.
Recent systems work on mixed isolation regimes, such as Milano et al.’s MixT
\cite{milano2020mixt}, addresses related concerns operationally; our contribution
is a semantic and explanatory account that makes such mixtures explicit and
checkable.

\paragraph{Homeostasis.}
Gehrke, Korn, and Srivastava’s work on \emph{homeostasis}~\cite{gehrke2002homeostasis}
studies queries whose results remain invariant under admissible input changes,
enabling incremental maintenance without recomputation.
In our frame, homeostatic queries inhabit the \emph{depth 0} fragment:
the specification is already resolved; additional semantic commitments are
neither required nor permitted.

Our framework is complementary.
Rather than enforcing invariance, we model and explain what happens when
invariance does \emph{not} hold—when specifications intentionally admit multiple
admissible outcomes, and semantic commitments must be staged to make meaning
determinate.
Interpretation depth characterizes the structure and necessity of those stages,
while robustness questions diagnose which conclusions persist despite the
absence of finer semantic commitments.
Viewed this way, homeostasis isolates a class of specifications whose semantics
are already resolved, so incremental change does not require introducing new
semantic commitments.

\subsection{Interpretation Complexity}

This work builds on a prior complexity theory of interpretation
introduced in~\cite{hellerstein2026interpretation}, which characterizes the cost,
depth, and coordination requirements of semantic resolution.

The present paper focuses instead on explanation.
Interpretations are assumed given, and interpretation provenance explains how
observations depend on them.
Connections to descriptive complexity, fixpoint alternation, oracle models, and
coordination are discussed in detail in~\cite{hellerstein2026interpretation}.

% ================================================================
\section{Conclusion}
\label{sec:conclusion}
% ================================================================

This paper reframes provenance around a simple observation: many modern systems do
not produce a single semantic outcome from a history, but a space of admissible
ones.
In such settings, explaining an observation requires accounting for how ambiguity
was resolved—not merely how derivations proceeded after the fact.
Interpretation provenance makes those semantic commitments explicit and
composable, recovering classical algebraic provenance once resolution is fixed
while retaining visibility into the assumptions that made resolution possible.

Our development shows that this perspective is not domain-specific.
In transactional systems, interpretation provenance explains how query
observations depend on ordering and isolation commitments.
In logic programming with negation, it exposes the semantic dependencies that
govern resolution.
Across domains, certificates provide a minimal interface for carrying forward
exactly the information required to support further reasoning.

Interpretation provenance also opens a broader research agenda.

\paragraph{Robustness and semantic sensitivity.}
Once observations are annotated with the interpretations under which they hold, it
becomes natural to ask which are robust across interpretations and which are
fragile.
This enables provenance questions about semantic stability, policy change, and
compatibility across regimes that lie beyond traditional why- and how-provenance.
Developing principled robustness measures—qualitative or quantitative—remains an
open direction.

\paragraph{Engineering interpretation provenance.}
While interpretations are minimal by construction, the choice of
\emph{interpretive predicate basis} is not.
Designing effective predicate bases—balancing parsimony against the desire to
expose semantic sensitivity—raises new engineering questions at the boundary of
specification and provenance.
Coarser bases yield more robust observations but collapse distinctions; finer
bases increase discriminating power but can amplify fragility and cost, as
illustrated by isolation-levels in transactional systems (Section~\ref{sec:robustness}).
Understanding how to design specifications and predicate bases that make the
right semantic distinctions explicit, while supporting compact certificates and
efficient maintenance, is a central direction for future work.

\paragraph{Beyond databases.}
Semantic ambiguity and resolution arise well beyond data management.
Consistency models, coordination protocols, access-control systems, and machine
learning pipelines all exhibit underspecification followed by structured
commitment.
Interpretation provenance provides a vocabulary for reasoning across these
settings, suggesting a unifying foundation for provenance in ambiguous
computational systems.

The success of algebraic provenance lay in identifying the right semantic object
to explain.
Interpretation provenance extends that insight: when meaning itself depends on
resolution, provenance must explain the assumptions that made meaning possible.

\appendix
\section{Appendix: Worked Examples}
\label{sec:appendix-examples}

\subsection{Full transactional provenance example}
\label{sec:appendix-concurrency}

We begin with a minimal example that already requires interpretation provenance.

\paragraph{Data, query, and history.}
Let $S(k,y)$ be a relation and consider the set-valued query $Q(y) : \exists k.\; S(k,y)$.
The database initially contains $S(0,a)$ with provenance variable $x_a$.
Five transactions execute:
\[
\begin{array}{ll}
T_1: \textsf{insert } S(1,b) \ (x_1) \qquad &
T_2: \textsf{insert } S(2,b) \ (x_2)\\
T_3: \textsf{insert } S(3,c) \ (x_3) \qquad &
T_4: \textsf{insert } S(4,d) \ (x_4)\\
T_5: \textsf{delete } S(4,d). &
\end{array}
\]
A query transaction $T_Q$ evaluates $Q$ under a
serializable execution model, where the visibility of updates is
determined by the chosen serialization order.

\paragraph{Semantic ambiguity.}
The history alone does not determine whether $d$ is visible to $T_Q$.
That depends on how $T_Q$ is ordered relative to $T_4$ and $T_5$.
Write $\varphi_{T_i \prec T_j}$ for the interpretive predicate committing
that the schedule that ensued corresponds to equivalent serial schedule in which
$T_i$ committed prior to $T_j$.

Transactions $T_1,T_2,T_3$ conflict with the query $Q$ and therefore must
be ordered with respect to $T_Q$ in any serializable execution. 
For expository purposes, we restrict attention to serializable executions
in which $T_Q$ is last in the serialization; this loses no generality for
the example, and allows the query result to serve as a witness of the final
database state.
Transactions $T_4$ and $T_5$ conflict with each other and with $T_Q$,
leaving their relative ordering unresolved.

Given that analysis, two resolving interpretations suffice:
\[
\begin{aligned}
  I_{\mathsf{in}} &\;\triangleq\; \varphi_{T_4 \prec T_Q} \wedge \varphi_{T_Q \prec T_5} &&\quad\text{(observe insert before delete)}\\
  I_{\mathsf{out}} &\;\triangleq\; \varphi_{T_5 \prec T_Q} &&\quad\text{(observe after delete).}
\end{aligned}
\]

\paragraph{Depth--0 provenance under each interpretation.}
Conditioned on $I_{\mathsf{in}}$, the resolved snapshot contains $b,c,d$:
\[
  \Prov_{Q \mid I_{\mathsf{in}}}(b)=x_1+x_2,\;
  \Prov_{Q \mid I_{\mathsf{in}}}(c)=x_3,\;
  \Prov_{Q \mid I_{\mathsf{in}}}(d)=x_4.
\]
Conditioned on $I_{\mathsf{out}}$, the snapshot contains $b,c$ but not $d$.

\paragraph{Why vs.\ why-not depends on interpretation.}
Under $I_{\mathsf{in}}$, the natural classical question for $d$ is \emph{why},
answered by $x_4$.
Under $I_{\mathsf{out}}$, the natural question is \emph{why-not} $d$: the witness
$k=4$ fails because $S(4,d)$ is false in the resolved snapshot.
There is no single classical provenance question for $d$ prior to fixing an
interpretation. Classical why-not provenance explains absence relative to a fixed resolved
semantics.
Under semantic ambiguity, absence admits additional structure; a brief discussion
appears in Appendix~\ref{app:whynot}.

\paragraph{Interpretation provenance.}
Interpretation provenance attaches the resolving interpretation to the ordinary
derivational explanation.
For the observed output $\{b,c\}$, we record
\[
  \mathsf{IP}_Q(b)\triangleq (I,\, x_1+x_2),\quad
  \mathsf{IP}_Q(c)\triangleq (I,\, x_3),
\]
where $I$ is the interpretation assumed by the system (e.g., $I_{\mathsf{out}}$).

This single representation supports both why- and why-not reasoning.
Under $I_{\mathsf{in}}$, $d$ admits a why-provenance explanation; under
$I_{\mathsf{out}}$, it admits a why-not explanation.
Interpretation provenance makes this dependence explicit.

\subsubsection{Certificates for Robustness Queries}

In this example, a certificate sufficient for classical why/why-not provenance of $d$ \emph{under a fixed serialization} can be as small as the relevant ordering commitments (i.e., the resolving interpretation assumed by the system, such as $I_{\mathsf{in}}$ or $I_{\mathsf{out}}$). Conditioned on that fixed interpretation, $d$ is either present (with why-provenance witness $x_4$) or absent (with a why-not explanation tied to the same fixed commitments).

By contrast, the robustness query “Does $d$ appear under all serializable executions?” requires the
\emph{strictly stronger} certificate needed for robustness: it must preserve the fact that the history does
\emph{not} impose a total order between the insert $T_4$ and delete $T_5$ relative to $T_Q$, so that both
$I_{\mathsf{in}}$ and $I_{\mathsf{out}}$ remain admissible under the specification prior to resolution. This
is directly analogous to the negation-cycle case in Appendix~\ref{sec:appendix-datalog} and illustrates the same
query-family effect discussed in Section~\ref{sec:certificates}.

\subsection{Datalog\texorpdfstring{$^{\neg}$}{ with negation} worked interpretation provenance}
\label{sec:appendix-datalog}

\subsubsection{Stratified semantics: explicit interpretation provenance}

\paragraph{Interpretation provenance (example).}
Write $\varphi_{\mathrm{stab}(p)}$ for the interpretive predicate that commits
that the lower stratum defining $p$ has stabilized for the current history, and
$\varphi_{\mathrm{neg}(q)}$ for the predicate that commits to the stratified
reading of $q \leftarrow \neg p$ \emph{given} that commitment.
A resolving interpretation for the stratified fragment can be written
\[
  I_{\mathrm{strat}} \;\triangleq\;
  \varphi_{\mathrm{stab}(p)} \wedge \varphi_{\mathrm{neg}(q)}.
\]

The interpretive predicate dependency is strict:
\[
  \varphi_{\mathrm{stab}(p)} \;\rightarrow\; \varphi_{\mathrm{neg}(q)}.
\]
The negation commitment cannot be applied until stabilization of the lower stratum
has been resolved.

Over the resolved specification $\Spec \mid I_{\mathrm{strat}}$, $p(c)$ is derived
from $a(c)$ and has ordinary derivation provenance
$k_p \triangleq \Prov(a(c))$.
Interpretation provenance records that this derivation is justified under the
stabilization-and-negation commitments:
\[
  \mathsf{IP}(p(c)) \;\triangleq\; (I_{\mathrm{strat}},\, k_p) .
\]

In contrast, $q(c)$ is determined only after $\varphi_{\mathrm{stab}(p)}$ is in
force; under our $H$ with $a(c)$, $q(c)$ does not hold under stratified semantics.

\subsubsection{Well-founded semantics: explicit interpretation provenance}

In our example, the stratified fragment induces the same prefix dependencies as
before:
\[
  \varphi_{\mathrm{stab}(p)} \;\rightarrow\; \varphi_{\mathrm{neg}(q)}.
\]
Beyond this prefix, the remaining atoms participate in a negation cycle.
Neither $r(c)$ nor $s(c)$ is justified without circular reasoning, so WFS commits both
to \emph{undefined}.

\paragraph{Interpretation provenance (example).}
Let $\mathbf{t},\mathbf{f},\mathbf{u}$ denote the well-founded truth values
true/false/undefined.
A resolving interpretation for $(P,H)$ under WFS can be written schematically as
\[
  I_{\mathrm{wfs}} \;\triangleq\;
  \varphi_{p(c)=\mathbf{t}} \wedge \varphi_{q(c)=\mathbf{f}}
  \wedge \varphi_{r(c)=\mathbf{u}} \wedge \varphi_{s(c)=\mathbf{u}} .
\]

The interpretive predicate dependency graph therefore consists of a stratified
prefix followed by a cycle:
\[
  \varphi_{\mathrm{stab}(p)} \;\rightarrow\; \varphi_{\mathrm{neg}(q)}
  \qquad\text{and}\qquad
  \varphi_{r(c)=\mathbf{u}} \;\leftrightarrow\; \varphi_{s(c)=\mathbf{u}} .
\]
The cycle admits no precedence ordering.
Well-founded semantics resolves it deterministically by committing \emph{both}
atoms to $\mathbf{u}$.

Over $\Spec \mid I_{\mathrm{wfs}}$, derivational provenance is depth--0; in
particular $p$ again has
$k_p \triangleq \Prov(a(c))$.
Interpretation provenance records the semantic commitments required for that
derivation:
\[
  \mathsf{IP}(p(c)) \;\triangleq\; (I_{\mathrm{wfs}},\, k_p) .
\]

The distinctive feature is that $I_{\mathrm{wfs}}$ resolves the negation cycle
without making a two-valued choice: it commits explicitly to assigning $r(c)$ and $s(c)$
the value $\mathbf{u}$.

\subsubsection{Stable semantics: explicit interpretation provenance}

As before, the stratified fragment induces the same prefix dependencies:
\[
  \varphi_{\mathrm{stab}(p)} \;\rightarrow\; \varphi_{\mathrm{neg}(q)}.
\]
Beyond this prefix, the negation cycle admits multiple two-valued resolutions.

For our example with $a(c)$, the stratified fragment forces $p(c)$ true and $q(c)$ false,
while the cycle admits two stable outcomes: one with $r(c)$ true and $s(c)$ false, and
one with $s(c)$ true and $r(c)$ false.

\paragraph{Interpretation provenance (example).}
Accordingly, stable semantics yields multiple resolving interpretations for the
same $(P,H)$.
Write
\[
  I_{\mathrm{stab}}^{(r)} \triangleq
  \varphi_{p(c)=\mathbf{t}} \wedge \varphi_{q(c)=\mathbf{f}}
  \wedge \varphi_{r(c)=\mathbf{t}} \wedge \varphi_{s(c)=\mathbf{f}}
\]
and
\[
  I_{\mathrm{stab}}^{(s)} \triangleq
  \varphi_{p(c)=\mathbf{t}} \wedge \varphi_{q(c)=\mathbf{f}}
  \wedge \varphi_{s(c)=\mathbf{t}} \wedge \varphi_{r(c)=\mathbf{f}} .
\]

The dependency structure is therefore:
\[
  \varphi_{\mathrm{stab}(p)} \;\rightarrow\; \varphi_{\mathrm{neg}(q)},
  \qquad
  \varphi_{r(c)=\mathbf{t}} \;\leftrightarrow\; \varphi_{s(c)=\mathbf{f}},
  \qquad
  \varphi_{s(c)=\mathbf{t}} \;\leftrightarrow\; \varphi_{r(c)=\mathbf{f}} .
\]
Unlike WFS, stable semantics does not prescribe a unique resolution of the cycle;
each choice yields a distinct resolved interpretation.

Within either resolved outcome, derivational provenance is ordinary:
$p(c)$ again has $k_p \triangleq \Prov(a(c))$, so
\[
  \mathsf{IP}(p(c)) \;\triangleq\; (I_{\mathrm{stab}}^{(r)},\, k_p)
  \quad\text{or}\quad
  \mathsf{IP}(p(c)) \;\triangleq\; (I_{\mathrm{stab}}^{(s)},\, k_p) ,
\]
depending on which stable interpretation is assumed.
The cycle atoms exhibit genuinely interpretation-dependent conclusions; e.g.,
\[
  \mathsf{IP}(r(c)) \;\triangleq\; (I_{\mathrm{stab}}^{(r)},\, 1)
  \qquad\text{and}\qquad
  \mathsf{IP}(s(c)) \;\triangleq\; (I_{\mathrm{stab}}^{(s)},\, 1) .
\]
Thus, in the stable setting, interpretation provenance must account explicitly for
which \emph{resolved interpretation} was selected.

\subsubsection{Semantic change example: negation semantics}

This example compares outcomes across specifications, not across interpretations of a single fixed specification.

This example is a concrete instance of inter-specification semantic change in the sense of
Section~\ref{sec:robustness}: holding the program and extensional history fixed,
different negation semantics induce different resolving interpretations, and thus
different conditioned outcomes.

\paragraph{Program and instance.}
Consider the program $P$, consisting of a stratified fragment defining $p(c)$,
and a negation cycle between $q(c)$ and $r(c)$:
\[
\begin{array}{l}
p(c) \;\drule\; a(c).\\
q(c) \;\drule\; \neg r(c).\\
r(c) \;\drule\; \neg q(c).
\end{array}
\]
with extensional instance $H = \{a(c)\}$.

\paragraph{Well-founded semantics.}
Under well-founded semantics, the negation cycle is unresolved.
The unique resolving interpretation commits:
\[
I_{\mathrm{wfs}}
\;\triangleq\;
\varphi_{p(c)=\mathbf{t}}
\;\wedge\;
\varphi_{q(c)=\mathbf{u}}
\;\wedge\;
\varphi_{r(c)=\mathbf{u}} .
\]
(These $\varphi$ are interpretive predicates at the meta-level: they constrain the
admissible outcomes induced by the program and history, rather than denoting
object-language Datalog atoms.)
Over $\Spec \mid I_{\mathrm{wfs}}$, $p(c)$ is true, while $q(c)$ and $r(c)$ are undefined.

\paragraph{Stable model semantics.}
Under stable model semantics, the same program admits two resolving interpretations:
\[
I^{(q)}_{\mathrm{stab}}
\;\triangleq\;
\varphi_{p(c)=\mathbf{t}}
\wedge \varphi_{q(c)=\mathbf{t}}
\wedge \varphi_{r(c)=\mathbf{f}},
\qquad
I^{(r)}_{\mathrm{stab}}
\;\triangleq\;
\varphi_{p(c)=\mathbf{t}}
\wedge \varphi_{r(c)=\mathbf{t}}
\wedge \varphi_{q(c)=\mathbf{f}} .
\]

These resolving interpretations are mutually incompatible conditionings of the
same underlying program and history.
In particular, $p(c)$ persists under both conditionings, while $q(c)$ and $r(c)$
depend on which resolving interpretation is assumed.

\subsubsection{Certificates and Semantic-Change Queries}

This appendix has emphasized certificates sufficient to support derivational provenance once an interpretation is fixed. To see how provenance query families enlarge certificate requirements (Section~\ref{sec:certificates}), consider the negation-cycle program used in the semantic-change examples:
\[
  q(c) \;\drule\; \neg r(c).\qquad
  r(c) \;\drule\; \neg q(c).
\]
Under a fixed resolving interpretation (e.g., committing $p(c)=\mathbf{t}$ in the surrounding example), a minimal certificate sufficient for the derivational provenance of $p(c)$ need only preserve the commitments that make $p(c)$ well-defined (and hence justify $k_p=\Prov(a(c))$). However, answering a semantic-change query such as “Would $q(c)$ become true under stable-model semantics?” requires a \emph{strictly stronger} certificate: it must preserve the \emph{unresolved negation structure} of the cycle (the mutual negative dependency between $q(c)$ and $r(c)$), not merely the resolved truth assignments under one semantics.

This is an instance of Theorem~6.1 (Query Monotonicity): supporting a richer query (semantic change across semantics/interpretations) can only increase the information a certificate must retain. This illustrates the query-parametric notion of certificate sufficiency formalized in Section~\ref{sec:certificates}.

% ================================================================
\section{Appendix: Interpretation Provenance for Systems Lineage}
\label{app:systems-lineage}
% ================================================================

This appendix illustrates interpretation provenance in a systems setting that
lies outside traditional database semantics.
The example is adapted directly from the Dapper distributed tracing model
introduced by Sigelman et al.~\cite{sigelman2010dapper}, abstracting its trace
semantics rather than its implementation details.

The example serves two purposes.
First, it shows that interpretation provenance extends the database tradition of
provenance to forms of \emph{systems lineage} that have historically resisted
formal treatment.
Second, it demonstrates how the parsimony of interpretations---as minimal semantic
commitments---can refine systems lineage explanations by making explicit which
assumptions are required to justify diagnostic conclusions.

\subsection{Trace Specification and Semantic Ambiguity}

A distributed trace consists of a set of \emph{spans}, each representing an
invocation of a service.
Each span $X$ has a client-side send and receive event
$(\mathsf{cs}_X,\mathsf{cr}_X)$ and a server-side receive and send event
$(\mathsf{sr}_X,\mathsf{ss}_X)$.
Client and server events occur on different hosts with unsynchronized clocks.

The trace specification $\Spec$ maps a history $H$ of timestamped span events to a
set of admissible executions $O$.
An execution assigns each span a concrete placement on a global timeline and
induces a partial order over spans.
The specification enforces only causal constraints:
\[
  \mathsf{cs}_X \prec \mathsf{sr}_X
  \qquad\text{and}\qquad
  \mathsf{ss}_X \prec \mathsf{cr}_X ,
\]
together with parent--child relationships recorded in the trace metadata.

Because clocks are unsynchronized, these constraints determine only \emph{bounds}
on server-side timestamps.
As a result, multiple executions---with different relative orderings and overlaps
of spans---are admissible under the same observed history.
Critical-path analysis is therefore semantically ambiguous prior to interpretation.

\subsection{Diagnostic Query: Critical Path}

We consider the diagnostic query:
\begin{quote}
\emph{Which span lies on the critical path of the request?}
\end{quote}
Formally, the query maps an execution to the span (or set of spans) whose duration
determines the end-to-end latency of the root span.

Under a fixed execution, this query is deterministic.
Under the specification $\Spec$, however, different admissible executions may
yield different critical paths.

\subsection{Interpretive Predicates}

Interpretation resolves this ambiguity by committing to semantic assumptions that
fix a particular execution.
Relevant interpretive predicates include:
\begin{itemize}
  \item $\varphi_{\mathrm{align}(h_i,h_j)}$, committing to a concrete clock
        alignment between hosts $h_i$ and $h_j$ within the admissible bounds;
  \item $\varphi_{\mathrm{order}(X,Y)}$, committing that span $X$ precedes span
        $Y$ when their relative order is otherwise unconstrained;
  \item $\varphi_{\mathrm{cp}=X}$, committing that $X$ is the critical-path
        witness once ordering and alignment are fixed.
\end{itemize}

These predicates are ambiguity-sensitive and may be semantically dependent.
For example, committing to $\varphi_{\mathrm{cp}=X}$ presupposes commitments that
fix the relative placement of competing spans.

\subsection{Resolving Interpretations}

Let $B$ and $C$ be sibling spans whose relative order is not determined by causal
constraints.
Two resolving interpretations suffice to illustrate the provenance structure:
\[
  I_B \;\triangleq\;
  \varphi_{\mathrm{align}} \wedge \varphi_{\mathrm{order}(B,C)} \wedge
  \varphi_{\mathrm{cp}=B},
\]
\[
  I_C \;\triangleq\;
  \varphi_{\mathrm{align}} \wedge \varphi_{\mathrm{order}(C,B)} \wedge
  \varphi_{\mathrm{cp}=C}.
\]
Both interpretations resolve $\Spec$, but they induce different critical-path
outcomes.

\subsection{Interpretation Provenance}

Once an interpretation is fixed, critical-path computation is ordinary.
Let $k_B$ (resp.\ $k_C$) denote the derivational justification that $B$ (resp.\ $C$)
is the longest-duration span in the resolved execution.
Interpretation provenance records:
\[
  \mathsf{IP}(\mathrm{cp}=B) \;\triangleq\; (I_B,\, k_B),
  \qquad
  \mathsf{IP}(\mathrm{cp}=C) \;\triangleq\; (I_C,\, k_C) .
\]

The provenance makes explicit that the diagnostic conclusion depends on semantic
commitments about clock alignment and span ordering.
Without fixing an interpretation, the question ``which span is on the critical
path?'' is ill-posed.

\subsection{Certificates and Parsimony}

Notably, the full trace history is unnecessary to justify the diagnosis.
A sufficient certificate consists of:
\begin{itemize}
  \item the parent--child span structure;
  \item the inequality constraints induced by send/receive events;
  \item the specific alignment and ordering commitments recorded in $I_B$ or
        $I_C$.
\end{itemize}

Interpretation provenance therefore supports a compressed form of systems lineage:
only the semantic commitments required for the diagnostic observation need be
retained.
This refines existing tracing practice by separating observable evidence from the
interpretive assumptions used to explain it.

\paragraph{Summary.}
This example demonstrates that interpretation provenance applies naturally to
systems diagnostics.
It extends provenance beyond data derivations to explanations of execution-level
properties, while retaining a minimal, explicit account of the semantic
assumptions required to justify those explanations.

\section{Absence Under Semantic Ambiguity}
\label{app:whynot}
In settings with semantic ambiguity, absence admits distinctions that cannot be
expressed at interpretation depth~0.

Given a specification $\Spec$ and history $H$, a conclusion may be absent in
several semantically distinct ways:
(i) it fails under a particular resolved interpretation;
(ii) it fails under all resolving interpretations of the fixed specification
$\Spec$; or
(iii) it holds under some interpretations of a \emph{different} specification
$\Spec'$—for example, under a different negation semantics or isolation level—but
under none of the interpretations admitted by $\Spec$ itself.
Classical why-not provenance collapses these cases, because it conditions on a
single resolved semantics and a fixed specification.

Interpretation provenance makes the first two cases explicit.
By recording the semantic commitments under which observations hold, it explains
absence in terms of the non-existence (or variability) of supporting
interpretations of $\Spec$, rather than in terms of blocked derivations.
The third case is an inter-specification question and is handled via semantic
change and preservation analysis (Section~\ref{sec:robustness}), rather than by
interpretation provenance alone.
\bibliographystyle{ACM-Reference-Format}
\bibliography{references}

\end{document}
